% GENERATED FILE. Edit canonical lessons, then run npm run generate:books.
% canonical-source-hash: fnv1a64:9b34c3d4ec025209
% canonical-lessons: RU-C437-zabluditsya, RU-C437-oriyentirovatsya, RU-C437-bluzhdat, RU-C437-kompas, RU-C437-navigator

\chapter{Getting lost}
\label{ch:ru-getting-lost}
\hlchaptermodality{437}

\begin{chapteropening}
\textbf{By the end of this chapter:} \emph{I can say I am lost: getting lost, finding my way, wandering, the compass and the satnav.}

Complete the last lesson of chapter 437: I can say I am lost: getting lost, finding my way, wandering, the compass and the satnav.
\end{chapteropening}

\section[\texorpdfstring{\ru{заблудиться} (zabludítsya)}{заблудиться (zabludítsya)}]{\ru{заблудиться} (zabludítsya) — to get lost}

\label{lesson:RU-C437-zabluditsya}



\begin{quote}
Before the new one: say the Russian for to fix, then the Russian for a receipt.
\end{quote}

\subsection*{You'll want to know: \ru{заблудиться}}
\textbf{\ru{заблудиться}} (\emph{zabludítsya}) — \textquotedblleft{}to get lost\textquotedblright{}.

\begin{quote}
\ru{В} \ru{старом} \ru{городе} \ru{мы} \ru{быстро} \ru{заблудились}.
\end{quote}

\begin{quote}
\emph{V stárom górode my býstro zabludílis.}
\end{quote}

\begin{quote}
We quickly got lost in the old town.
\end{quote}

Perfective. Future \emph{\ru{заблужусь}, \ru{заблудишься}}.

\begin{cousinweb}
\emph{\ru{заблудиться}} is built from \emph{\ru{за}-}, \emph{\ru{блудить}} 'to wander, to stray' and \emph{-\ru{ся}}; EtymDB lists it with the sense 'to get lost'. The root is Proto-Slavic \emph{blǫdŭ} (EtymDB), the same root as in \emph{\ru{блуждать}} 'to wander'.
\end{cousinweb}

\begin{grammarlens}[title={What you've built}]
One more word for telling, explaining and describing at B1.
\end{grammarlens}

\subsection*{Your turn}
\begin{itemize}
\raggedright
  \item \emph{Say it:} \emph{zabludítsya}
  \item \emph{Say it:} \emph{zabludítsya}, once more
  \item \emph{Say it:} a sentence of your own with \emph{zabludítsya}
  \item \emph{Recall:} say the Russian for to fix, then the Russian for a receipt, then say \emph{zabludítsya} again
\end{itemize}

\begin{tcolorbox}[breakable,colback=teal!4,colframe=teal!35!black,title={Before you move on}]
What does \ru{заблудиться} mean? (\textquotedblleft{}To get lost\textquotedblright{}.) Say it once more.
\end{tcolorbox}
\section[\texorpdfstring{\ru{ориентироваться} (oriyentírovatsya)}{ориентироваться (oriyentírovatsya)}]{\ru{ориентироваться} (oriyentírovatsya) — to find one's way, to get one's bearings}

\label{lesson:RU-C437-oriyentirovatsya}



\begin{quote}
Before the new one: say the Russian for a receipt, then the Russian for to get lost.
\end{quote}

\subsection*{You'll want to know: \ru{ориентироваться}}
\textbf{\ru{ориентироваться}} (\emph{oriyentírovatsya}) — \textquotedblleft{}to find one's way, to get one's bearings\textquotedblright{}.

\begin{quote}
\ru{В} \ru{новом} \ru{городе} \ru{я} \ru{плохо} \ru{ориентируюсь}.
\end{quote}

\begin{quote}
\emph{V nóvom górode ya plókho oriyentíruyus.}
\end{quote}

\begin{quote}
I am bad at finding my way in a new town.
\end{quote}

Imperfective; perfective \emph{\ru{сориентироваться}}. Takes \emph{\ru{в}} + prepositional.

\begin{cousinweb}
\emph{\ru{ориентироваться}} is usually traced to French \emph{orienter} 'to orient, to set facing north' (EtymDB), on Latin \emph{oriēns} 'rising; the east' (Lewis \& Short).
\end{cousinweb}

\begin{grammarlens}[title={What you've built}]
One more word for telling, explaining and describing at B1.
\end{grammarlens}

\subsection*{Your turn}
\begin{itemize}
\raggedright
  \item \emph{Say it:} \emph{oriyentírovatsya}
  \item \emph{Say it:} \emph{oriyentírovatsya}, once more
  \item \emph{Say it:} a sentence of your own with \emph{oriyentírovatsya}
  \item \emph{Recall:} say the Russian for a receipt, then the Russian for to get lost, then say \emph{oriyentírovatsya} again
\end{itemize}

\begin{tcolorbox}[breakable,colback=teal!4,colframe=teal!35!black,title={Before you move on}]
What does \ru{ориентироваться} mean? (\textquotedblleft{}To find one's way, to get one's bearings\textquotedblright{}.) Say it once more.
\end{tcolorbox}
\section[\texorpdfstring{\ru{блуждать} (bluzhdát)}{блуждать (bluzhdát)}]{\ru{блуждать} (bluzhdát) — to roam, to wander}

\label{lesson:RU-C437-bluzhdat}



\begin{quote}
Before the new one: say the Russian for to get lost, then the Russian for to find one's way.
\end{quote}

\subsection*{You'll want to know: \ru{блуждать}}
\textbf{\ru{блуждать}} (\emph{bluzhdát}) — \textquotedblleft{}to roam, to wander\textquotedblright{}.

\begin{quote}
\ru{Мы} \ru{два} \ru{часа} \ru{блуждали} \ru{по} \ru{лесу}.
\end{quote}

\begin{quote}
\emph{My dva chasá bluzhdáli po lésu.}
\end{quote}

\begin{quote}
We wandered around the forest for two hours.
\end{quote}

Imperfective; takes \emph{\ru{по}} + dative. Present \emph{\ru{блуждаю}}.

\begin{cousinweb}
\emph{\ru{блуждать}} is built on \emph{\ru{блуд}} (EtymDB), from Proto-Slavic \emph{blǫdŭ}; the same root gives \emph{\ru{заблудиться}} 'to get lost'.
\end{cousinweb}

\begin{grammarlens}[title={What you've built}]
One more word for telling, explaining and describing at B1.
\end{grammarlens}

\subsection*{Your turn}
\begin{itemize}
\raggedright
  \item \emph{Say it:} \emph{bluzhdát}
  \item \emph{Say it:} \emph{bluzhdát}, once more
  \item \emph{Say it:} a sentence of your own with \emph{bluzhdát}
  \item \emph{Recall:} say the Russian for to get lost, then the Russian for to find one's way, then say \emph{bluzhdát} again
\end{itemize}

\begin{tcolorbox}[breakable,colback=teal!4,colframe=teal!35!black,title={Before you move on}]
What does \ru{блуждать} mean? (\textquotedblleft{}To roam, to wander\textquotedblright{}.) Say it once more.
\end{tcolorbox}
\section[\texorpdfstring{\ru{компас} (kómpas)}{компас (kómpas)}]{\ru{компас} (kómpas) — a compass}

\label{lesson:RU-C437-kompas}



\begin{quote}
Before the new one: say the Russian for to find one's way, then the Russian for to roam.
\end{quote}

\subsection*{You'll want to know: \ru{компас}}
\textbf{\ru{компас}} (\emph{kómpas}) — \textquotedblleft{}a compass\textquotedblright{}.

\begin{quote}
\ru{Без} \ru{компаса} \ru{в} \ru{лесу} \ru{легко} \ru{заблудиться}.
\end{quote}

\begin{quote}
\emph{Bez kómpasa v lesú legkó zabludítsya.}
\end{quote}

\begin{quote}
Without a compass it is easy to get lost in the forest.
\end{quote}

Masculine. Sailors stress the second syllable: \emph{kompás}.

\begin{cousinweb}
\emph{\ru{компас}} was borrowed from German \emph{Kompass} (EtymDB).
\end{cousinweb}

\begin{grammarlens}[title={What you've built}]
One more word for telling, explaining and describing at B1.
\end{grammarlens}

\subsection*{Your turn}
\begin{itemize}
\raggedright
  \item \emph{Say it:} \emph{kómpas}
  \item \emph{Say it:} \emph{kómpas}, once more
  \item \emph{Say it:} a sentence of your own with \emph{kómpas}
  \item \emph{Recall:} say the Russian for to find one's way, then the Russian for to roam, then say \emph{kómpas} again
\end{itemize}

\begin{tcolorbox}[breakable,colback=teal!4,colframe=teal!35!black,title={Before you move on}]
What does \ru{компас} mean? (\textquotedblleft{}A compass\textquotedblright{}.) Say it once more.
\end{tcolorbox}
\section[\texorpdfstring{\ru{навигатор} (navigátor)}{навигатор (navigátor)}]{\ru{навигатор} (navigátor) — a satnav, a navigator}

\label{lesson:RU-C437-navigator}



\begin{quote}
Before the new one: say the Russian for to roam, then the Russian for a compass.
\end{quote}

\subsection*{You'll want to know: \ru{навигатор}}
\textbf{\ru{навигатор}} (\emph{navigátor}) — \textquotedblleft{}a satnav, a navigator\textquotedblright{}.

\begin{quote}
\ru{Навигатор} \ru{показал} \ru{нам} \ru{дорогу}.
\end{quote}

\begin{quote}
\emph{Navigátor pokazál nam dorógu.}
\end{quote}

\begin{quote}
The satnav showed us the way.
\end{quote}

Masculine. Also a ship's or a plane's navigator.

\begin{cousinweb}
\emph{\ru{навигатор}} is usually traced to Latin \emph{nāvigātor} 'a sailor' (EtymDB lists the Latin word but not the Russian), from \emph{nāvigāre} 'to sail', on \emph{nāvis} 'a ship' (Lewis \& Short).
\end{cousinweb}

\begin{grammarlens}[title={What you've built}]
One more word for telling, explaining and describing at B1.
\end{grammarlens}

\subsection*{Your turn}
\begin{itemize}
\raggedright
  \item \emph{Say it:} \emph{navigátor}
  \item \emph{Say it:} \emph{navigátor}, once more
  \item \emph{Say it:} a sentence of your own with \emph{navigátor}
  \item \emph{Recall:} say the Russian for to roam, then the Russian for a compass, then say \emph{navigátor} again
\end{itemize}

\begin{tcolorbox}[breakable,colback=teal!4,colframe=teal!35!black,title={Before you move on}]
What does \ru{навигатор} mean? (\textquotedblleft{}A satnav, a navigator\textquotedblright{}.) Say it once more.
\end{tcolorbox}
